% Extended reading and original technical analysis. Included after the six core pages.
\par
\section{Nền tảng ra quyết định: trạng thái, MDP và quan sát}
\subsection{Từ một chương trình phản ứng đến bài toán quyết định tuần tự}
Một hành động có thể thay đổi những lựa chọn còn lại trong tương lai. Ví dụ, robot đi
vào một hành lang có thể tiết kiệm thời gian trước mắt nhưng làm mất cơ hội chọn đường
khác nếu cửa phía sau đóng. Vì vậy, việc đánh giá hành động cần xét chuỗi kết quả,
không chỉ phản hồi ngay lập tức. Tổng quan của Kaelbling, Littman và Moore trình bày
học tăng cường dưới góc nhìn tương tác thử--sai và lý thuyết quyết định Markov
\cite{rlsurvey}. Phần dưới cụ thể hóa các khái niệm bằng mô hình của báo cáo.

Một Markov Decision Process có thể biểu diễn bằng $(\mathcal S,\mathcal A,P,R,\gamma)$:
$\mathcal S$ là tập trạng thái, $\mathcal A$ là tập hành động, $P(s'|s,a)$ là xác suất
chuyển trạng thái, $R$ là phần thưởng và $\gamma$ là hệ số chiết khấu. Chính sách
$\pi(a|s)$ mô tả cách chọn hành động. Tính Markov đòi hỏi trạng thái hiện tại chứa đủ
thông tin để mô tả phân phối bước tiếp theo khi đã biết hành động.

\subsection{Tại sao cách định nghĩa trạng thái quan trọng?}
Trong GridWorld tĩnh của thực nghiệm, vị trí $(r,c)$ là trạng thái đủ: cùng một hành
động từ cùng vị trí tạo ra cùng kết quả. Nhưng trong ví dụ cạnh B--D đóng sau bước
đầu, vị trí đơn lẻ chưa mô tả đầy đủ môi trường. Để mô hình đầy đủ là Markov, có thể
thêm trạng thái đóng/mở cửa hoặc biến thời gian liên quan. Agent lại chỉ quan sát
được một phần thông tin, nên mô hình nội bộ là ước lượng về môi trường.

Ví dụ tự xây dựng: hai lịch sử đều đưa robot về B nhưng một lịch sử đã có thông báo
B--D bị đóng. Nếu chỉ lưu B, robot có thể chọn lại hành động lỗi. Khi trạng thái nội
bộ được bổ sung tập cạnh chặn, hai tình huống được phân biệt. Tuy nhiên, bộ nhớ đó vẫn
không đại diện hoàn hảo cho thế giới nếu cửa có thể tự mở lại mà robot không biết.

\subsection{Quan sát không đầy đủ và trạng thái niềm tin}
Khi cảm biến không cho biết trạng thái thật, có thể duy trì một phân phối niềm tin
$b_t(s)$ thay cho một kết luận chắc chắn \cite{pomdp}. Theo phép cập nhật Bayes, một dạng cập nhật là
\[
b_{t+1}(s')=\eta\,O(o_{t+1}|s',a_t)\sum_s P(s'|s,a_t)b_t(s),
\]
trong đó $O$ mô tả xác suất nhận quan sát, còn $\eta$ chuẩn hóa tổng xác suất bằng 1.
Đây là công thức nền để hiểu vấn đề, không phải phần đã cài trong notebook. Ví dụ,
cảm biến cửa báo sai 10\% thì một lần báo ``đóng'' không nên được hiểu như bằng chứng
chắc chắn vĩnh viễn. Ngược lại, công cụ mô phỏng ở đây trả về lỗi xác định nên dùng
tập cạnh chặn đơn giản là phù hợp với giả định thực nghiệm.

\begin{center}\small
\begin{tabularx}{\linewidth}{@{}p{3.3cm}XX@{}}\toprule
\textbf{Môi trường} & \textbf{Thông tin cần giữ} & \textbf{Hệ quả thiết kế}\\\midrule
Tĩnh, quan sát đầy đủ & Trạng thái hiện tại và mô hình & Có thể tìm kế hoạch trước\\
Động, quan sát cục bộ & Lịch sử liên quan và thời điểm & Cần kiểm tra lại giả định\\
Có nhiễu & Độ tin cậy hoặc phân phối niềm tin & Tránh quyết định chắc chắn quá sớm\\
Đối tác cũng hành động & Trạng thái phối hợp và cam kết & Cần giải quyết xung đột\\\bottomrule
\end{tabularx}\end{center}
\textbf{Liên hệ với LLM agent:} lịch sử hội thoại là một nguồn quan sát, không tự động
là trạng thái đúng của nghiệp vụ. Một câu ``đã thanh toán'' trong tin nhắn và trạng
thái xác nhận từ hệ thống thanh toán phải được phân biệt về nguồn và độ tin cậy.

\par
\section{LLM agent, workflow và hợp đồng sử dụng công cụ}
\subsection{Phân biệt phạm vi của hai cách dùng thuật ngữ agent}
Trong cách phân loại AI cổ điển, bộ tìm đường có mục tiêu là một agent. Trong bài viết
kỹ thuật về hệ thống LLM, Anthropic phân biệt workflow có đường điều khiển định trước
với agent để LLM tự định hướng quá trình và lựa chọn công cụ \cite{effective}.
Hai cách dùng có phạm vi khác nhau. Báo cáo sử dụng nghĩa rộng cho agent tìm kiếm và
agent RL, đồng thời ghi rõ chính sách của trợ lý tra cứu là mã Python xác định.

LLM agent có thể được xem như bộ đề xuất hành động làm việc cùng một bộ thực thi.
Mô hình nhận mô tả mục tiêu, công cụ và trạng thái rồi đề xuất tên tool cùng tham số.
Bộ thực thi kiểm tra đề xuất, gọi tool, lưu quan sát và quyết định bước tiếp theo.
Nếu một lời gọi không hợp lệ, việc trả lỗi có cấu trúc giúp phân biệt lỗi lập kế
hoạch với lỗi hạ tầng. Đặt quyền kiểm soát ở bộ thực thi giúp hành vi của hệ thống
có thể kiểm tra mà không phụ thuộc hoàn toàn vào cách mô hình diễn đạt.

\subsection{Tool use có thể được học như thế nào?}
Toolformer nghiên cứu việc huấn luyện mô hình chọn API, thời điểm gọi, tham số và cách
kết hợp kết quả API vào dự đoán văn bản \cite{toolformer}. Đây là một hướng học sử
dụng công cụ; nó khác với việc lập trình sẵn một hàm chọn tool hoặc chỉ liệt kê tool
trong prompt. Khả năng gọi đúng API vẫn cần được đánh giá trên đầu vào và điều kiện
thực thi cụ thể.

Trong thiết kế minh họa của báo cáo, công cụ \texttt{search} chỉ nhận truy vấn và
\texttt{top\_k}; \texttt{read} chỉ nhận ID tài liệu. Hợp đồng này xác định rõ miền
đầu vào, cấu trúc đầu ra và lỗi. Mã không cho phép tùy ý chuyển một chuỗi thành lệnh
hệ điều hành. Ngay cả khi thay bộ chọn hành động bằng LLM, lớp hợp đồng này vẫn cần
được giữ và kiểm tra độc lập.

\subsection{Phân biệt các loại lỗi để xử lý đúng}
\begin{center}\small
\begin{tabularx}{\linewidth}{@{}p{3.0cm}XX@{}}\toprule
\textbf{Loại lỗi} & \textbf{Ví dụ tự xây dựng} & \textbf{Cách xử lý đề xuất}\\\midrule
Schema & \texttt{top\_k=true} & Từ chối kiểu sai; yêu cầu sửa tham số\\
Không có dữ liệu & Không tìm thấy ID & Không tạo ra tài liệu giả\\
Tạm thời & Timeout dịch vụ đọc & Retry có giới hạn, ghi mã yêu cầu\\
Quyền hạn & Không được đọc hồ sơ & Dừng hoặc chuyển người có quyền\\
Mâu thuẫn trạng thái & Phòng vừa được đặt & Đọc lại và lập kế hoạch khác\\\bottomrule
\end{tabularx}\end{center}

\subsection{Hành động đọc, hành động ghi và kiểm chứng}
Một truy vấn tra cứu không tạo ra cùng trách nhiệm như gửi thư hoặc đặt chỗ. Với hành
động ghi, hệ thống cần xác định người chịu trách nhiệm, điều kiện cho phép và khả năng
khôi phục. Ví dụ tự thiết kế: chuẩn bị một bản nháp lịch học có thể tự động; ghi lịch
lên tài khoản dùng chung cần tuân theo quy trình xác nhận của đơn vị. Tool trả về
``accepted'' cũng có thể chỉ cho biết yêu cầu đã xếp hàng, chưa hoàn tất.

Trong notebook, kết quả cuối là trạng thái vị trí hoặc văn bản có ID nguồn, nên có
phép kiểm tra trực tiếp. Đối với hệ thống thật, nên tách trạng thái đã yêu cầu, đang
thực hiện, đã xác nhận và thất bại. Cấu trúc trạng thái này giúp agent không báo hoàn
thành chỉ vì đã tạo được một lời gọi hợp lệ.

\par
\section{Bộ nhớ, ngữ cảnh, reflection và thư viện kỹ năng}
\subsection{Bộ nhớ có nhiều chức năng}
Tổng quan của Wang và cộng sự mô tả kiến trúc LLM agent theo các mô-đun như hồ sơ,
bộ nhớ, lập kế hoạch và hành động \cite{llmsurvey}. Để áp dụng, báo cáo phân biệt
những dữ liệu mà agent cần giữ theo mục đích sử dụng, thay vì xem mọi văn bản đã lưu
là cùng một loại tri thức.

\begin{center}\small
\begin{tabularx}{\linewidth}{@{}p{2.7cm}XX@{}}\toprule
\textbf{Loại thông tin} & \textbf{Ví dụ của báo cáo} & \textbf{Câu hỏi kiểm tra}\\\midrule
Ngữ cảnh làm việc & Mục tiêu, vị trí, bước đang chạy & Còn phù hợp lượt hiện tại không?\\
Ký ức tình huống & B--D bị chặn sau lần thử & Nguồn nào xác nhận? Khi nào hết hạn?\\
Tri thức tham khảo & Tài liệu quy trình có phiên bản & Có quyền dùng và còn hiệu lực không?\\
Kỹ năng thực thi & Hàm kiểm tra lịch trống & Đã kiểm thử với tool phiên bản nào?\\
Sở thích người dùng & Ưu tiên phòng gần thư viện & Có đồng ý lưu lâu dài không?\\\bottomrule
\end{tabularx}\end{center}

\subsection{Reflection thay đổi cái gì?}
Reflexion dùng phản hồi bằng ngôn ngữ và bộ nhớ tình huống để cải thiện các lần thử,
không cần cập nhật trọng số mô hình \cite{reflexion}. Điều này giúp phân biệt học
trong ngữ cảnh với tối ưu tham số. Từ ``reinforcement'' trong tên phương pháp không
có nghĩa mọi thành phần của nó giống thuật toán Q-learning dạng bảng.

Ví dụ tự phân tích: sau một lần tạo phiếu sai đơn vị tiếp nhận, agent lưu nhận xét
``thiết bị mạng thuộc nhóm hạ tầng, không thuộc nhóm máy tính cá nhân''. Trước nhiệm
vụ tiếp theo, nó truy hồi nhận xét đó. Nếu quy trình tổ chức thay đổi, ghi nhớ cũ có
thể dẫn đến sai mới. Vì vậy, phản hồi được lưu nên kèm phạm vi áp dụng, thời điểm và
nguồn xác nhận, đồng thời có cơ chế sửa hoặc quên.

Generative Agents nghiên cứu việc lưu trải nghiệm bằng ngôn ngữ, tổng hợp reflection
và truy hồi để lập hành vi trong một môi trường mô phỏng xã hội \cite{generative}.
Kết quả hành vi có vẻ hợp lý trong mô phỏng không tự động là bằng chứng về độ chính
xác của một hệ thống xử lý hồ sơ. Thước đo cần phù hợp với loại nhiệm vụ.

\subsection{Từ ghi nhớ một câu trả lời đến tái sử dụng một kỹ năng}
Voyager kết hợp chương trình học tự động, thư viện kỹ năng dạng code và cải tiến dựa
trên phản hồi môi trường trong Minecraft \cite{voyager}. Ý tưởng đáng chú ý là tái
sử dụng hành vi có thể thực thi, thay vì chỉ lưu đoạn hội thoại. Việc lấy một kỹ năng
từ thư viện vẫn đòi hỏi điều kiện đầu vào phù hợp và kiểm tra kết quả.

Trong mở rộng do báo cáo đề xuất, kỹ năng ``tìm phòng phù hợp'' có thể đóng gói thành
hàm nhận ngày, giờ, số người, rồi trả về danh sách phòng. Hàm này nên được phiên bản
hóa, kiểm thử trường hợp không có phòng và không tự ý đặt chỗ. Khi dùng lại, agent
cần kiểm tra lược đồ lịch hiện tại; nếu trường dữ liệu đã đổi, việc tái sử dụng code
cũ có thể làm mất lợi ích của bộ nhớ kỹ năng.

\subsection{Chi phí và giới hạn của việc lưu tất cả}
Đưa toàn bộ lịch sử vào mỗi lần gọi mô hình làm ngữ cảnh dài hơn và chứa cả thông tin
không còn liên quan. Tóm tắt giúp giảm độ dài nhưng có thể bỏ mất điều kiện quan trọng.
Một giải pháp thiết kế là lưu bản ghi gốc, tạo bản tóm tắt có tham chiếu và chỉ truy
hồi những phần liên quan đến quyết định hiện tại. Đánh giá bộ nhớ cần đo cả độ đúng
khi truy hồi, dữ liệu lỗi thời và mức giảm chi phí sau khi tính công sức duy trì.

\par
\section{RAG, truy hồi thông tin và dẫn chứng}
\subsection{RAG giải quyết nhu cầu kiến thức bên ngoài mô hình}
Lewis và cộng sự giới thiệu cách kết hợp bộ nhớ tham số của mô hình sinh với bộ nhớ
phi tham số được truy hồi từ kho tài liệu \cite{rag}. Trong mô hình RAG-Sequence,
một cách biểu diễn là:
\[
p(y|x)\approx\sum_{z\in\mathrm{TopK}(x)}p_{\eta}(z|x)\,p_{\theta}(y|x,z),
\]
với $x$ là câu hỏi, $z$ là tài liệu truy hồi, $y$ là câu trả lời; $\eta$ và $\theta$
lần lượt biểu diễn thành phần truy hồi và sinh. Công thức thể hiện rằng tìm tài liệu
và sử dụng tài liệu là hai bài toán riêng, cùng ảnh hưởng kết quả.

\subsection{Luồng dữ liệu để triển khai một hệ thống có nguồn}
Một thiết kế tham khảo của báo cáo gồm: thu nhận tài liệu được phép dùng; gán ID và
phiên bản; tách đoạn; lập chỉ mục; truy hồi; chọn bằng chứng; tạo câu trả lời; kiểm tra
liên kết giữa phát biểu và nguồn. Mỗi công đoạn có thể làm sai thông tin. Chẳng hạn,
tách đoạn làm mất tiêu đề ``chỉ áp dụng cho học kỳ hè'' có thể khiến một mốc thời gian
được áp dụng sai cho cả năm học.

Trong bài thực hành 2, kho chỉ có bốn quy định giả lập. Để dễ đọc toàn bộ code, dùng
vector đếm token và cosine:
\[
\operatorname{sim}(q,d)=\frac{\sum_w q_w d_w}{\sqrt{\sum_w q_w^2}\sqrt{\sum_w d_w^2}}.
\]
Nếu một vector rỗng, mã trả score 0. Đây là truy hồi lexical; nó không nhận ra mọi
quan hệ đồng nghĩa và không phải dense retrieval của nghiên cứu RAG. Câu trả lời được
trích nguyên từ tài liệu giả lập kèm ID, không gọi mô hình sinh nên không gán nhãn đây
là một bản tái lập RAG đầy đủ.

\subsection{Agentic retrieval khác một truy vấn duy nhất ở đâu?}
Trong một thiết kế mở rộng, agent có thể nhận ra bằng chứng thiếu, sửa câu hỏi, chọn
nguồn khác hoặc tra thêm tài liệu. Tuy nhiên, một chuỗi cố định search--read--answer
vẫn là workflow nếu mọi bước đã được mã hóa. Bài thực hành dùng workflow này để làm
rõ hợp đồng công cụ và nguồn gốc dữ liệu, không tuyên bố LLM tự điều phối.

Ví dụ tự xây dựng: câu hỏi ``được mượn bao nhiêu sách và trong bao lâu?'' cần hai dữ
kiện. Một đoạn nêu ``3 sách'' mà thiếu thời hạn chưa đủ trả lời. Hệ thống nên kiểm tra
mức bao phủ từng phần câu hỏi, thay vì chọn đoạn có score cao nhất rồi kết luận đã
đủ bằng chứng. Nếu không tìm được thời hạn, câu trả lời cần nêu phần chưa xác nhận.

\subsection{Dẫn chứng đúng gồm nhiều điều kiện}
\begin{center}\small
\begin{tabularx}{\linewidth}{@{}p{3.0cm}X@{}}\toprule
\textbf{Điều kiện} & \textbf{Phép kiểm tra đề xuất}\\\midrule
Nguồn tồn tại & ID/URL mở được, tài liệu còn truy cập được\\
Đúng phạm vi & Đúng đơn vị, thời gian, đối tượng, phiên bản\\
Hỗ trợ phát biểu & Đoạn trích thực sự chứng minh câu được viết\\
Không bỏ điều kiện & Giữ ngoại lệ và giới hạn ảnh hưởng kết luận\\
Đủ bao phủ & Mỗi phát biểu cần chứng cứ đều có nguồn phù hợp\\\bottomrule
\end{tabularx}\end{center}
Một citation có hình thức đẹp không bù được suy luận vượt quá tài liệu. Trong báo cáo
này, số liệu của thực nghiệm gắn với code và output; thông tin nghiên cứu gắn với bài
gốc; các ví dụ nghiệp vụ được ghi là thiết kế riêng. Cách tách nguồn này giúp người
đọc biết điều gì đã được kiểm tra và điều gì là hướng đề xuất.

\par
\section{Các chiến lược lập kế hoạch và tìm kiếm}
\subsection{Tìm kiếm trên mô hình tường minh}
BFS phù hợp khi mọi cạnh có chi phí bằng nhau. Nếu chi phí khác nhau nhưng không âm,
có thể dùng Dijkstra. A* sắp xếp ứng viên theo $f(n)=g(n)+h(n)$, với $g(n)$ là chi phí
đã đi và $h(n)$ là ước lượng chi phí còn lại \cite{astar}. Điều kiện của heuristic và quy tắc xử lý
đỉnh đã xét ảnh hưởng bảo đảm tối ưu; không nên gán bảo đảm của tìm kiếm cổ điển cho
một hàm đánh giá bằng LLM mà chưa kiểm chứng.

Ví dụ tự tính: đường A--B--F có hai cạnh với chi phí 1 và 9, tổng 10; đường A--C--D--F
có ba cạnh với chi phí 2, 2 và 2, tổng 6. BFS tối ưu số cạnh nên có thể chọn đường đầu;
mục tiêu tối ưu chi phí yêu cầu đường sau. Vì vậy, kết luận ``đường ngắn nhất'' luôn
phải chỉ rõ ngắn theo số bước, độ dài vật lý hay thời gian dự kiến.

\subsection{Lập kế hoạch trong không gian ngôn ngữ}
Tree of Thoughts nghiên cứu việc xét nhiều hướng suy luận, đánh giá và quay lui ở
cấp các bước lập luận bằng văn bản \cite{tot}. Cách tiếp cận này mở rộng không gian
tìm kiếm nhưng cũng tạo nhu cầu đánh giá các ứng viên. Một bước lập luận có vẻ hợp
lý chưa chắc thỏa mãn ràng buộc của bài toán gốc.

Giả sử tác vụ là lập lịch ba môn vào hai phòng. Agent sinh ba lịch ứng viên; bộ kiểm
tra xác định việc trùng thời gian, sức chứa và thiết bị bắt buộc. LLM có thể giúp đề
xuất lịch, còn các ràng buộc kiểm được bằng chương trình nên được kiểm bằng chương
trình. Khi không lịch nào hợp lệ, có thể mở rộng tìm kiếm hoặc báo không đủ tài nguyên,
thay vì chọn lịch có lời giải thích thuyết phục nhất.

\subsection{Lập trước và điều chỉnh theo quan sát}
Theo cách tổ chức ReAct, kết quả hành động tham gia quyết định bước kế tiếp
\cite{react}. ReWOO tách phần lập kế hoạch với việc thu thập kết quả công cụ
\cite{rewoo}. Trong phân tích của báo cáo, hai nhu cầu cạnh tranh là tránh lập lại
những phần đã ổn định và tránh bám vào kế hoạch đã mất hiệu lực.

\begin{center}\small
\begin{tabularx}{\linewidth}{@{}p{3.0cm}XX@{}}\toprule
\textbf{Tình huống} & \textbf{Chiến lược đề xuất} & \textbf{Phép kiểm tra}\\\midrule
Các truy vấn độc lập & Chuẩn bị và gom truy vấn & Kết quả có cùng thời điểm không?\\
Phụ thuộc kết quả trước & Quan sát rồi chọn bước tiếp & Còn thiếu dữ kiện nào?\\
Không gian chọn lớn & Sinh ứng viên và kiểm ràng buộc & Ứng viên có thực thi được không?\\
Môi trường biến động & Kế hoạch ngắn và điểm kiểm tra & Giả định nào vừa thay đổi?\\\bottomrule
\end{tabularx}\end{center}

\subsection{Phân rã nhiệm vụ cần giữ điều kiện toàn cục}
Nếu chia một chuyến đi thành đặt vé, đặt phòng và sắp xếp lịch, từng tác vụ có thể
thành công riêng nhưng tổng thể vẫn thất bại khi giờ nhận phòng không khớp giờ đến.
Do đó, bộ điều phối cần giữ các ràng buộc liên kết, thứ tự phụ thuộc và tiêu chí hoàn
thành chung. Kế hoạch nên là một cấu trúc có trạng thái, thay vì danh sách câu văn
không thể biết bước nào đã xác nhận.

Trong notebook tìm đường, chuỗi đỉnh là kế hoạch thực thi cụ thể. Khi một cạnh lỗi,
agent bỏ phần kế hoạch còn lại và tìm từ vị trí hiện tại. Đây là một minh họa rõ về
việc kế hoạch phải được kiểm tra với môi trường, dù cơ chế suy luận không dùng LLM.

\par
\section{Hệ đa tác tử: giao tiếp, điều phối và trách nhiệm}
\subsection{Vì sao dùng nhiều tác tử?}
AutoGen nghiên cứu xây dựng ứng dụng bằng các agent có thể hội thoại, kết hợp mô hình,
tool và đầu vào con người \cite{autogen}. Một giá trị thiết kế là tách vai trò,
nguồn dữ liệu hoặc quyền thao tác. Việc tăng số agent cũng tăng số điểm phối hợp cần
kiểm tra; không thể suy ra hiệu quả chỉ từ số lượng tác tử.

Ví dụ tự đề xuất: hệ thống hỗ trợ báo cáo có tác tử tìm tài liệu, tác tử viết và tác tử
kiểm chứng. Tác tử tìm tài liệu trả các bản ghi ID, tiêu đề, đoạn hỗ trợ; tác tử viết
chỉ dùng dữ kiện đã có; tác tử kiểm chứng đối chiếu từng phát biểu với nguồn. Nếu cả
ba chỉ lặp lại cùng một thông tin sai từ một tài liệu, sự đồng thuận không tạo thêm
chứng cứ. Đầu ra cần được kiểm tra theo nguồn và điều kiện nhiệm vụ.

\subsection{Các mô hình tổ chức có thể lựa chọn}
\begin{center}\small
\begin{tabularx}{\linewidth}{@{}p{2.7cm}XX@{}}\toprule
\textbf{Cấu trúc} & \textbf{Cách hoạt động} & \textbf{Điểm cần quản lý}\\\midrule
Chuỗi vai trò & Nghiên cứu $\to$ soạn $\to$ kiểm tra & Lỗi đầu chuỗi lan xuống cuối\\
Điều phối trung tâm & Một bộ phân việc cho chuyên gia & Nút điều phối có thể thành điểm nghẽn\\
Làm song song & Các tác vụ độc lập cùng chạy & Cần hợp nhất và xử lý mâu thuẫn\\
Phản biện & Một bên đề xuất, một bên kiểm & Giới hạn vòng tranh luận và chuẩn chấm\\\bottomrule
\end{tabularx}\end{center}

Trong cấu trúc đề xuất của báo cáo, thông điệp nên có mã nhiệm vụ, mục tiêu, nguồn dữ
liệu, trạng thái, kết quả và lỗi. Một yêu cầu ``kiểm tra lại'' thiếu tiêu chí sẽ khiến
tác tử nhận không biết phải đối chiếu điều gì. Nếu đầu ra là bản vá code, tiêu chí có
thể là bài kiểm tra tái hiện lỗi và tập kiểm tra hồi quy; nếu đầu ra là tổng hợp tài
liệu, tiêu chí có thể là đúng trích dẫn và đủ câu hỏi được trả lời.

\subsection{Mô hình chi phí và lỗi phối hợp}
Với các phần việc tuần tự có thời gian $T_1,\dots,T_n$, tổng độ trễ gần
$\sum_i T_i$. Khi phần việc độc lập được chạy đồng thời, độ trễ lý tưởng gần
$\max_i T_i$, nhưng còn thời gian điều phối, tổng hợp và chờ tài nguyên. Tổng lượng
tính toán có thể tăng dù người dùng nhận kết quả sớm hơn.

Một hệ đa tác tử cũng có thể gặp tranh chấp khi hai agent cùng sửa một hồ sơ. Đề xuất
thiết kế là ghi phiên bản dữ liệu đã đọc và từ chối cập nhật nếu phiên bản thay đổi.
Với thông điệp bất đồng bộ, cần phân biệt nhiệm vụ đang chờ, đang chạy, đã hoàn thành
và đã hủy, đồng thời quy định tác tử nào có quyền ra quyết định cuối.

\subsection{Kiểm tra lợi ích bằng đối chứng}
Để chứng minh mô hình đa tác tử hữu ích, có thể so với một agent dùng cùng công cụ,
cùng dữ liệu và tổng ngân sách tương đương. Các tiêu chí là hoàn thành, chi phí,
độ trễ và lỗi. Nếu mô hình nhiều tác tử được cấp số lượt gọi gấp mười lần, kết quả
tốt hơn chưa cho biết lợi ích đến từ tổ chức vai trò hay chỉ từ nhiều tài nguyên hơn.

Trong phạm vi thực hành này, không cần đa tác tử để giải đồ thị sáu đỉnh hoặc kho bốn
tài liệu. Phần đa tác tử là kiến thức mở rộng và phân tích thiết kế; báo cáo không
tuyên bố đã đo hiệu quả của một framework đa tác tử trong các thực nghiệm đi kèm.

\par
\section{Học tăng cường: Q-learning và thiết kế thí nghiệm}
\subsection{Giá trị hành động và mục tiêu cập nhật}
Watkins và Dayan trình bày Q-learning như phương pháp cải thiện ước lượng chất lượng
hành động trong môi trường Markov; kết quả hội tụ cần các điều kiện thích hợp
\cite{qlearning}. Công thức được triển khai trong thực nghiệm là
\[
Q(s,a)\leftarrow Q(s,a)+\alpha\left[r+\gamma(1-d)\max_{a'}Q(s',a')-Q(s,a)\right],
\]
trong đó $d=1$ khi đạt trạng thái kết thúc. $\alpha$ là learning rate, $\gamma$ là mức
chiết khấu. Tại terminal, không cộng lợi ích của một hành động tương lai không xảy ra.

Ví dụ tính tay để đối chiếu code: $Q(s,a)=0$, phần thưởng đến đích là 10 và
$\alpha=0.2$. Vì trạng thái tiếp theo là terminal, giá trị mới bằng
$0+0.2(10-0)=2$. Cho dù hàng Q của đích chứa số 100 để thử lỗi, phép cập nhật vẫn phải
ra 2. Cell kiểm tra cố tình tạo tình huống này để phát hiện bootstrap sai.

\subsection{Khám phá và khai thác}
Chính sách $\epsilon$-greedy chọn hành động ngẫu nhiên với xác suất $\epsilon$, còn lại
chọn hành động có Q cao nhất. Trong code, $\epsilon$ giảm từ khoảng 0.9 theo hệ số
0.995 mỗi episode và có sàn 0.05. Tie-break ngẫu nhiên có seed tránh luôn ưu tiên
hướng đầu danh sách khi các giá trị Q bằng nhau.

Khám phá gây thêm bước và có thể làm reward huấn luyện thấp hơn reward đánh giá của
chính sách đóng băng. Hai đại lượng này không nên trộn trong cùng kết luận.
Sau huấn luyện, việc đánh giá đặt $\epsilon=0$ và giữ nguyên Q. Baseline ngẫu nhiên
dùng cùng lưới, vị trí đầu và giới hạn bước để so sánh hành vi thực tế.

\subsection{Các lựa chọn cụ thể trong bài thực hành}
\begin{center}\small
\begin{tabularx}{\linewidth}{@{}p{3.6cm}X@{}}\toprule
\textbf{Thành phần} & \textbf{Thiết lập thực nghiệm của báo cáo}\\\midrule
Không gian trạng thái & Lưới $5\times5$; 5 ô tường; đích $(4,4)$\\
Hành động & Lên, phải, xuống, trái; hành động sai giữ nguyên vị trí\\
Reward & Đến đích $+10$; bước thường $-1$; hành động sai $-3$\\
Huấn luyện & 2.500 episode/seed; seed 7, 42, 2026; tối đa 100 bước/episode\\
Siêu tham số & $\alpha=0.2$, $\gamma=0.95$; Q khởi tạo bằng 0\\
Đánh giá & 100 episode/chính sách; không cập nhật Q; baseline ngẫu nhiên\\\bottomrule
\end{tabularx}\end{center}

\subsection{Termination, truncation và giới hạn kết luận}
Đến đích là terminal của bài toán. Hết 100 bước là dừng do giới hạn tính toán; mã vẫn
bootstrap theo trạng thái tiếp theo ở lần cập nhật cuối vì giới hạn này không được
định nghĩa là terminal trong MDP. Nếu thay đổi bài toán thành ``phải đến trước hạn'',
cần đưa thời gian còn lại vào trạng thái hoặc định nghĩa lại quá trình và hàm thưởng.

Học với $\alpha$ cố định trong hữu hạn episode và một bản đồ không chứng minh các điều
kiện hội tụ tổng quát. Ba seed làm rõ mức phụ thuộc vào khởi tạo ngẫu nhiên trong
thiết lập đã chọn; chúng chưa cho phép kết luận về khả năng chuyển sang bản đồ mới.
Đánh giá tổng quát cần thêm nhiều bản đồ, vị trí đầu khác nhau, nhiễu và phân tích
độ nhạy siêu tham số.

\par
\section{Đánh giá agent: benchmark, đối chứng và bằng chứng}
\subsection{Đánh giá tương tác khác đánh giá một câu trả lời}
AgentBench tổ chức đánh giá năng lực agent qua nhiều môi trường tương tác
\cite{agentbench}. WebArena cung cấp môi trường web có thể tái lập và chấm mức hoàn
thành chức năng của tác vụ \cite{webarena}. SWE-bench đánh giá sửa lỗi phần mềm dựa
trên yêu cầu và kho mã thực tế \cite{swebench}. Các bài này gợi ý rằng kết quả cuối
cần được kiểm tra trong môi trường, thay vì chỉ chấm câu giải thích của mô hình.

Các benchmark nói trên khác nhau về dữ liệu, quyền truy cập, công cụ và chuẩn chấm.
Báo cáo không so sánh trực tiếp điểm số giữa chúng và không dùng số liệu của các mô
hình cũ để mô tả hiệu năng hiện tại. Phần thực nghiệm ở đây là kiểm tra minh họa trên
đồ thị và dữ liệu nhỏ, không được gọi là đạt một benchmark nghiên cứu.

\subsection{Một quy trình đánh giá có thể áp dụng}
Trước hết, cố định mục tiêu và điều kiện thành công. Sau đó xác định tập tình huống,
ngân sách, seed và dữ liệu mà agent được nhìn thấy. Tiếp theo chạy đối chứng, lưu
nhật ký và phân loại lỗi. Cuối cùng đối chiếu kết quả với những giả định ban đầu.
Nếu thay đổi prompt hoặc siêu tham số sau khi xem toàn bộ tập kiểm tra, cần một tập
đánh giá khác để tránh kết quả bị tối ưu theo ví dụ đã biết.

\begin{center}\small
\begin{tabularx}{\linewidth}{@{}p{3.0cm}XX@{}}\toprule
\textbf{Khía cạnh} & \textbf{Phép đo đề xuất} & \textbf{Điều phải ghi kèm}\\\midrule
Đúng nhiệm vụ & Tỷ lệ đạt toàn bộ ràng buộc & Cách chấm và tình huống thất bại\\
Hiệu quả & Chi phí, độ trễ, số tool call & Cùng ngân sách với đối chứng\\
Ổn định & Phân bố kết quả theo seed & Độ ngẫu nhiên nằm ở thành phần nào\\
Khả năng phục hồi & Kết quả khi tool lỗi hoặc thiếu dữ liệu & Mức lỗi và chính sách retry\\
Nguồn dẫn & Độ bao phủ, phù hợp của citation & Phát biểu nào nguồn hỗ trợ\\\bottomrule
\end{tabularx}\end{center}

\subsection{Ablation để biết thành phần nào tạo ích lợi}
Với agent tìm đường, có thể bỏ bộ nhớ cạnh chặn để quan sát hành vi lặp lại lỗi; với
trợ lý tra cứu, có thể tắt cache để đo số lần gọi thêm; với Q-learning, có thể so với
chính sách ngẫu nhiên và kiểm tra ảnh hưởng lịch giảm $\epsilon$. Một ablation hợp
lý chỉ thay thành phần đang kiểm tra, giữ điều kiện khác tương đương.

Trong notebook, cache được kiểm bằng hai câu hỏi giống nhau: lần đầu có search và
read; lần sau trả kết quả có phiên bản từ cache và không gọi tool. Phép thử này chứng
minh hành vi cache trong ví dụ. Nó chưa chứng minh cache luôn an toàn khi tài liệu
thay đổi, khi quyền người dùng khác nhau hoặc khi câu hỏi có ngữ nghĩa tương tự nhưng
không giống chuỗi ký tự.

\subsection{Đọc tỷ lệ thành công đúng mức}
Trong GridWorld xác định, 100 episode đánh giá chính sách tốt trên cùng bản đồ không
phải 100 môi trường độc lập. Một kết quả 100/100 có thể chỉ phản ánh cùng đường đi
được lặp lại. Với các phép thử Bernoulli độc lập có tính đại diện, sai số chuẩn gần
$\sqrt{\hat p(1-\hat p)/n}$ có thể giúp mô tả bất định ở vùng thích hợp; công thức
xấp xỉ này không xử lý tốt biên 0 hoặc 1 và không sửa được thiên lệch chọn mẫu.
Vì thế, báo cáo ưu tiên ghi rõ tập dữ liệu và quy trình trước khi diễn giải một tỷ lệ.

\par
\section{Độ tin cậy, bảo mật và vòng đời vận hành}
\subsection{Dữ liệu bên ngoài có thể chứa chỉ dẫn không đáng tin}
Greshake và cộng sự nghiên cứu indirect prompt injection: chỉ dẫn có thể được cài vào
nội dung mà ứng dụng LLM truy hồi hoặc xử lý \cite{injection}. Hệ thống cần giữ ranh
giới giữa yêu cầu được ủy quyền và dữ liệu bên ngoài. Một trang web tự ghi ``hãy bỏ
qua quy định'' không trở thành quy tắc điều khiển ứng dụng chỉ vì agent đọc được nó.

Trong kiến trúc đề xuất, dữ liệu tài liệu được truyền như bằng chứng; bộ thực thi giữ
quyền quyết định công cụ nào được phép gọi. Việc chỉ cho phép \texttt{search} và
\texttt{read} trong bài thực hành hạn chế bề mặt hành động, nhưng không được diễn giải
thành chứng minh chống mọi tấn công. Không có LLM trong code để thử hành vi làm theo
chỉ dẫn độc hại; các kiểm tra hiện tại chỉ xác nhận hợp đồng công cụ đã cài.

\subsection{Ba lớp bảo vệ gắn với hành động}
Lớp thứ nhất kiểm tra đầu vào: kiểu dữ liệu, kích thước, ID và quyền truy cập. Lớp
thứ hai kiểm soát thực thi: giới hạn tool, timeout, ngân sách, cách ly và nhật ký.
Lớp thứ ba kiểm tra kết quả: xác nhận trạng thái, đối chiếu ràng buộc và quyết định
cần người can thiệp. Đây là phân rã thiết kế của báo cáo; hiệu quả thực tế phải được
kiểm thử theo ứng dụng cụ thể.

Ví dụ một agent tạo phiếu sửa thiết bị có thể được cho quyền đọc danh mục công khai
và ghi phiếu nháp, nhưng không được sửa danh sách người dùng. Nếu dữ liệu yêu cầu chứa
mã lệnh hoặc một câu đề nghị đổi quyền, bộ thực thi vẫn xử lý như nội dung yêu cầu.
Quyền ghi cần được giới hạn theo thao tác và đối tượng, thay vì một cờ ``được dùng API''
chung cho mọi loại cập nhật.

\subsection{Từ notebook sang một dịch vụ có thể vận hành}
\begin{center}\small
\begin{tabularx}{\linewidth}{@{}p{3cm}X@{}}\toprule
\textbf{Giai đoạn} & \textbf{Sản phẩm cần có trong thiết kế đề xuất}\\\midrule
Định nghĩa bài toán & Mục tiêu, người dùng, ràng buộc, điều kiện dừng\\
Chuẩn bị dữ liệu & Nguồn, quyền sử dụng, phiên bản, quy tắc cập nhật\\
Xây công cụ & Schema, xác thực, lỗi, giới hạn, kiểm tra lặp thao tác\\
Thử nghiệm & Tình huống chuẩn, ca biên, dữ liệu lỗi, đối chứng\\
Triển khai hạn chế & Giám sát, khả năng dừng, phương án chuyển người\\
Duy trì & Phân tích log, sửa nguồn, cập nhật tool, kiểm tra hồi quy\\\bottomrule
\end{tabularx}\end{center}

Nhật ký nên ghi hành động và quan sát có ích cho việc kiểm tra, đồng thời tránh lưu
dữ liệu cá nhân không cần thiết. Với hệ thống LLM, có thể ghi mã yêu cầu, tool, tham
số đã được lọc, mã lỗi, trạng thái cuối và phiên bản cấu hình. Không cần xem mọi câu
lập luận được sinh ra là bằng chứng chính xác về cơ chế bên trong; bằng chứng quan
trọng hơn là yêu cầu nào đã được thực thi và môi trường xác nhận điều gì.

\subsection{Kết nối các kiến thức mở rộng}
Agent là lớp ra quyết định trong một hệ thống lớn hơn. Tìm kiếm giải quyết việc chọn
chuỗi hành động từ mô hình; RAG cung cấp kiến thức có thể cập nhật; bộ nhớ giữ kinh
nghiệm; RL cải thiện hành vi từ tương tác; đa tác tử tổ chức các vai trò; đánh giá và
kiểm soát xác nhận hành vi đáng tin trong phạm vi đã định nghĩa. Những thành phần này
có thể kết hợp nhưng không thay thế nhau. Ba bài thực hành sau minh họa từng phần đủ
rõ để đọc hết code, chạy lại và đối chiếu output với nhận định trong báo cáo.
